\chapter{Перелёт в клетку}\label{ch:peretyot}

\textbf{Что вы узнете:} как послать дрона в клетку и как связаны расстояние, время и скорость.

\section{Команда code}

\kod{d.code(cell, seconds)} посылает дрона в центр клетки \kod{cell} по прямой за \kod{seconds} секунд.

\solfile{l06_move.py}

После взлёта идёт \kod{d.code(13, 1)}. Это \emph{стартовая клетка}: дрон должен стоять на клетке 13 и узнать там своё место, а симулятор считает, что дрон начинает под первой клеткой своей программы. Потом дрон перелетает в клетку 16: три клетки вправо, то есть 120\cm.

\begin{figure}[H]
  \centering
  % Лента l06_move.py: 5 + 5 + 1 + 4 + 3 = 18 секунд
\begin{tikzpicture}[x=0.8cm, font=\sffamily\small]
  \def\h{0.9}
  \foreach \a/\b/\t/\f in {0/5/{calibrate}/paper, 5/10/{takeoff}/accentsoft, 10/11/{}/paper, 11/15/{code(16, 4)}/accentsoft, 15/18/{land}/paper}{
    \draw[fill=\f, draw=ink!60, thick] (\a,0) rectangle (\b,\h);
    \node[font=\scriptsize, text=ink] at ({(\a+\b)/2},{\h/2}) {\t};
  }
  \node[font=\tiny, text=ink, rotate=90] at (10.5,0.45) {code(13,1)};
  \draw[-{Stealth}, ink] (0,-0.5) -- (18.9,-0.5) node[right] {$t$, с};
  \foreach \t in {0,5,10,11,15,18}{\draw[ink] (\t,-0.42) -- (\t,-0.58); \node[below, font=\scriptsize] at (\t,-0.58) {\t};}
\end{tikzpicture}

  \caption{Лента шоу \texttt{l06\_move.py}: $5 + 5 + 1 + 4 + 3 = 18$ секунд.}\label{fig:time}
\end{figure}

\section{Скорость выбираете вы}

Скорость дрона вы не задаёте. Она получается из расстояния и секунд:
\[
  \text{скорость} = \frac{\text{расстояние}}{\text{секунды}}.
\]
Расстояние от клетки 13 до клетки 16 равно $3\times 40 = 120$\cm. За 4 секунды получается $120/4 = 30$\cm/с. Если написать \kod{d.code(16, 1)}, скорость будет 120\cm/с, и проверка скажет:

\begin{center}
\small
«Дрон 1: переход на клетку 16 за 1 с (120 см, 120 см/с) быстрее порога 100 см/с: дрон может перелететь клетку»
\end{center}

Порог такой: \textbf{не быстрее 100\cm/с}. Если превысить, настоящий дрон может проскочить мимо клетки.

\begin{zametka}
Расстояние считается по прямой между центрами клеток. По диагонали оно длиннее: от клетки 13 до клетки 52 это $\sqrt{120^2 + 120^2}\approx 170$\cm, значит нужно не меньше 1,7 секунды.
\end{zametka}

\begin{probuy}
  \item Перелетите из клетки 13 в клетку 15 за 3 секунды. Чему равна скорость?
  \item Файл \texttt{bad\_examples/bad\_fast.py} (внизу) получает замечание. Исправьте его.
  \item Туда и обратно: из клетки 13 в клетку 17 и обратно, по 4 секунды в каждую сторону. Сколько секунд длится шоу?
  \item Перелетите из клетки 13 в клетку 52 по диагонали. За какое самое короткое время можно, чтобы не было замечания?
\end{probuy}
\solfile{bad_examples/bad_fast.py}
\otvety{6}
